\section{Decisions When the Evaluated Future Changes}\label{sec:r21future}
The query-volume experiment in Section~\ref{sec:r19experiment} keeps the
future reference policy fixed. We now change that future in the same capital
economy. The question is whether a learned continuation continues to support
accurate decisions, when a check requires refreshing it, and which economic
comparative statics can be certified from the returned actions.

\subsection{Economic design and frozen service rules}
The current postdecision map, temporary charge, utility weight, adjustment
term, and feasible withdrawal interval are held fixed across five future
regimes. Relative to the anchor, subsequent withdrawal changes from $0.55$
to $0.25$ or $0.85$, future production coupling changes from $0.08$ to $0.24$,
or the weight on subsequent and terminal payoffs doubles. The current
production coefficient remains $0.08$ in every regime. Thus the experiment
changes future policy, technology, and valuation, not merely the labels on
current tasks. The future policy remains specified in each regime; its
endogenous reoptimization is not assumed.

Each economy retains the sixteen-path stored innovation law, four dates, and
continuous common scalar action interval of the preceding experiment. State
dimensions are ten and fifty. There are $q\in\{1,16,256,1024\}$ prespecified task arguments per future and three new fixed fitting/cache
streams, giving 168 complete services and 840 future-specific outcomes.
``New'' distinguishes these declared streams from the preceding experiment;
it does not assert random sampling from an estimated economic population.
The task arguments and service code were frozen before the execution. The
same tasks are used across procedures, futures, and fitting streams within
each dimension--volume cell. Caches are reconstructed across streams.

Seven services compare reuse-only NBO, NBO that checks and refreshes, NBO
refitted for each future, check-and-refresh quadratic and radial-basis
continuations, cached SAA, and full-support enumeration. All inherit the same
fitting and action-selection routines. A reuse service first constructs the
anchor at the lower fitting budget. The adaptive rule checks each changed
future and, on failure, fits its continuation at the higher budget. The
fresh-refit rule tries the lower and then, if necessary, the higher budget
for each future. A service returns after its first successful check or
records that its budget was exhausted. It does not switch retrospectively
to the cheapest procedure. Every attempted fit, action query, verification,
and charged record write is retained.

The primary target is mean foregone transformed welfare at most $10^{-4}$
relative to the \emph{full scalar interval} in each future. The stopping
certificate evaluates the true finite-law objective and its curvature;
it does not use the surrogate's Hessian as if it were the objective's.
This certificate is deterministic, so no statistical confidence budget is
spent by repeated checks. A service that exhausts its budget without a
certificate is retained, even if its observed work is small. Diagnostic
prediction risks are computed only after the service clock is closed and
do not alter stopping.

\subsection{Accuracy, refresh, and actual work}
\begin{table}[htbp]\centering\small
\caption{Accuracy after changes in future economic primitives}\label{tab:r21accuracy}
\begin{tabular}{lrrrr}\toprule
Procedure & Complete & Futures & Refreshes & Largest bound\\ \midrule
NBO, reuse only & 0/24 & 72/120 & 0 & $2.71\times10^{-3}$\\
NBO, check and refresh & 24/24 & 120/120 & 48 & $5.29\times10^{-6}$\\
NBO, refit each future & 24/24 & 120/120 & 0 & $1.22\times10^{-5}$\\
Quadratic, check and refresh & 24/24 & 120/120 & 48 & $5.44\times10^{-5}$\\
RBF, check and refresh & 24/24 & 120/120 & 48 & $3.57\times10^{-6}$\\
Cached SAA & 24/24 & 120/120 & 0 & $3.66\times10^{-6}$\\
Full-support enumeration & 24/24 & 120/120 & 0 & $3.35\times10^{-18}$\\
\bottomrule\end{tabular}
\par\smallskip\begin{minipage}{.98\textwidth}\footnotesize A complete service contains all five futures. A successful future is certified against the full scalar interval at mean loss at most $10^{-4}$. A failed certificate is not a proof that the actual loss exceeds the tolerance. There are no execution exceptions. These are the original frozen outcomes, not new observations.\end{minipage}
\end{table}

NBO with checking and refresh certifies all 120 future-specific outcomes.
Its largest final mean loss bound is $5.288\times10^{-6}$, below the target
by more than an order of magnitude. Refit NBO likewise certifies every
future. Reuse-only NBO certifies the anchor and the two changes in future
withdrawal, but does not certify either the production or the future-weight
change in any of the 24 dimension--volume--stream cells. All 48 attempted
neural refreshes succeed. These results directly support the economic
use of a learned continuation under changing futures, while identifying
when treating an old evaluated object as current is inadequate.

The conventional adaptive surrogates, cached SAA, and enumeration also
certify every future. Successful neural decisions therefore do not imply
that their construction is uniquely necessary. Nor does the failure of a
certificate prove that actual loss exceeds its tolerance. The complete
record retains both distinctions.

\begin{table}[htbp]\centering\small
\caption{Recorded work for a five-future service (seconds)}\label{tab:r21cost}
\begin{tabular}{rrrrrrrrr}\toprule
$d$ & $q$ & Reuse & Adapt & Refit & Quad & RBF & SAA & Enum\\ \midrule
10 & 1 & 0.691 & 1.038 & 0.679 & 0.379 & 0.405 & 0.334 & 0.336\\
10 & 16 & 0.464 & 1.191 & 0.788 & 0.533 & 0.547 & 0.449 & 0.463\\
10 & 256 & 2.253 & 3.622 & 2.531 & 2.922 & 2.936 & 2.198 & 2.695\\
10 & 1,024 & 8.226 & 11.879 & 8.410 & 10.993 & 12.338 & 8.184 & 10.238\\
50 & 1 & 0.405 & 1.203 & 0.757 & 0.485 & 0.493 & 0.388 & 0.390\\
50 & 16 & 0.906 & 1.903 & 1.252 & 1.160 & 1.183 & 0.890 & 0.957\\
50 & 256 & 9.142 & 13.389 & 9.456 & 12.504 & 12.842 & 9.443 & 11.635\\
50 & 1,024 & 41.609 & 59.204 & 42.660 & 57.913 & 61.906 & 42.574 & 50.124\\
\bottomrule\end{tabular}
\par\smallskip\begin{minipage}{.98\textwidth}\footnotesize Columns follow Table~\ref{tab:r21accuracy}. Each cell is the arithmetic mean of three actually executed services, with $5q$ decisions per service. Costs include initialization, anchor fitting when used, all attempted fitting, queries, failed and successful checks, and charged durable writes. Reuse-only has unresolved futures and is not an accuracy-qualified winner. Common startup, post-stop diagnostics and the complete comparison remain separately recorded; no clock is replaced by a fitted line.\end{minipage}
\end{table}

Table~\ref{tab:r21cost} measures the amortization axis using complete
executed services rather than allocated prefixes. In every displayed cell,
checking and refreshing NBO costs more than refitting NBO for each future.
The former pays for its anchor and failed reuse checks and then, when
required, a higher-budget fit. The latter succeeds at its lower budget in
all five futures. This is an economically relevant cost of the registered
refresh rule, not a reason to omit its failures from the frontier. Cached
SAA and conventional surrogates remain serious alternatives; there is no
observed uniform neural cost advantage or inferred break-even outside the
measured query volumes.

The original full comparative clock is 1,897.326 seconds. Accounted deployed
services sum to 1,633.478 seconds and post-stop diagnostics to 256.996 seconds;
the remainder contains common startup and unallocated comparison work,
including metadata recording. The deployed clocks charge all stage checks
and failed attempts, but they are not complete process-isolation benchmarks.
The full comparison bill is reported rather than distributed into fictitious
method-specific observations. Its costs cannot be reduced after execution
by selecting a favorable row.

\subsection{Which continuation errors matter for the decision?}
The diagnostic action set is fixed exogenously at the lower endpoint,
midpoint, and upper endpoint of each feasible interval. All predictors are
assessed on that common set and, separately, at their own selected actions.
The targets are exact finite-law continuation means. Absolute risk,
action-centered risk, and decision loss are recorded separately, with no
Monte Carlo noise term silently absorbed into the target.

At 1,024 tasks, reusing the anchor after lowering future withdrawal gives
absolute continuation risk about $0.07084$ in dimension ten and $0.07111$
in dimension fifty. Yet centered risks are only $3.53\times10^{-8}$ and
$8.01\times10^{-8}$, and three-action mean decision losses are
$5.64\times10^{-7}$ and $2.50\times10^{-7}$ or smaller. A large continuation
level shift can therefore coexist with accurate action comparisons. Under
higher future production, the same unchanged neural continuation instead
has three-action losses about $5.32\times10^{-4}$ and $5.04\times10^{-4}$;
the registered check requires refreshing it. The common-action diagnostic
explains these outcomes through action-relevant errors rather than a ranking
of raw squared residuals. It is not itself the full-interval stopping
certificate. The supplement reports every procedure and future, including
the cases in which non-neural predictors are more accurate.

\subsection{An economic conclusion certified from NBO actions}
A resource charge reduces current withdrawals in the strictly concave
one-step objective, as derived in the preceding capital analysis. Changing
the future also changes the current optimum. We use each returned NBO
action's derivative and curvature certificate to enclose that optimum,
not merely its value. Section~\ref{app:r21proofs} proves
$|a-a^*|\leq r(a)/\kappa$ and the resulting paired interval calculation.

\begin{table}[htbp]\centering\small
\caption{Economic contrasts certified from NBO decisions}\label{tab:r21nbocontrasts}
\begin{tabular}{rlc}\toprule
$d$ & Future change & Mean optimal withdrawal minus anchor\\ \midrule
10 & Production & $[-0.035358,-0.027942]$\\
10 & Payoff weight & $[-0.018301,-0.013184]$\\
50 & Production & $[-0.037119,-0.026553]$\\
50 & Payoff weight & $[-0.019708,-0.012126]$\\
\bottomrule\end{tabular}
\par\smallskip\begin{minipage}{.98\textwidth}\footnotesize Each band uses only the final NBO actions and their true-objective derivative/curvature certificates. It encloses the true optimum response over the common 1,024 tasks. The printed band is the envelope of the three stream-specific bands, not an inferential interval over training randomness. All four upper endpoints are negative.\end{minipage}
\end{table}

For the fixed 1,024 tasks, the bands in Table~\ref{tab:r21nbocontrasts}
establish that stronger future production and a higher future-payoff weight
both lower mean optimal current withdrawal in each dimension. For example,
the fifty-dimensional production effect lies between approximately $-0.03712$
and $-0.02655$ in the model's withdrawal units. These are bands for the true
finite-law optimal response obtained from NBO's computed decisions. They
are not bands for the training algorithm's average action over an unobserved
stream population. Full-support enumeration supplies a sharper independent
benchmark in Table~\ref{tab:r21economic}; it confirms the signs rather than
being substituted for the neural certificate.

The exercise is a counterfactual comparative-static calculation in the
paper's specified capital economy, not an empirical estimate of an observed
reform. The welfare tolerance is in the transformed objective's units; the
external log-utility compensation account in the preceding section gives
its economic interpretation task by task. A bound on \emph{mean} value loss
must not be exponentiated and presented as a bound on mean compensation.
The numerical method supports a certified economic conclusion here, while
its comparative cost remains an empirical question answered by the full
service table.

\paragraph{Chronology and scope.}
The future-change source and streams belong to the frozen R20 execution.
This revision integrates that completed evidence, independently replays its
arithmetic certificates, and derives the displayed NBO response bands from
saved actions. The new reporting and interval contrasts are deterministic
post-execution analyses; they are not a newly randomized or prospectively
selected favorable experiment. All original 168 service records, 840 outcomes,
failed checks, source identities, and clocks remain available. The theoretical
composition result specifies the additional conditions for a full dynamic
policy certificate; the finite current-decision experiment is not substituted
for an uncomputed uniform or continuous-time error account.
